\section{Conclusiones y trabajo futuro}\label{i-0}

\subsection{Logros frente a los objetivos}\label{i-1}

Economicon reúne en un mismo prototipo lo que el proyecto se propuso: ingiere y normaliza el gasto de Azure, calcula indicadores, detecta anomalías y responde preguntas en lenguaje natural citando sus fuentes. Las cifras de coste las calcula siempre el servidor y nunca el modelo, y cada cliente solo ve sus datos. El sistema se construye, se prueba y se despliega de forma automatizada, y cada medición se puede repetir con un comando y consultar en Grafana. \cite{arquitectura-final,panel-evaluacion}

El alcance demostrado usa datos públicos y sintéticos y una API que simula Azure. Varias capacidades ---la recuperación híbrida, el chat generativo y el ciclo en AWS--- están implementadas pero apagadas por defecto o sin ejecutar en la nube, y su uso con una organización real está por demostrar. \cite{arquitectura-final,aws-ensayo}

\subsection{Lecciones y límites}\label{i-2}

La principal lección es que superar las pruebas técnicas no basta para aceptar un asistente. El chat cita bien, no inventa cifras y cuesta poco, pero acierta menos de lo exigido y tarda varias veces más que el objetivo (apartado~\ref{h-0}). Separar las cifras deterministas de la redacción del modelo resultó una buena decisión: los fallos se concentran en la redacción y en la verificación, no en los datos.

La segunda lección es de método. Medir con reglas, con una revisora y con un juez LLM da lecturas distintas que no se pueden sumar, y cambiar varias cosas a la vez impidió aislar el efecto del modelo. En el tramo final, además, varias integraciones se hicieron sin revisión cruzada para llegar a la entrega, y esas partes tienen menos verificación independiente. \cite{reevaluacion-modelo,continuidad}

\subsection{Trabajo futuro}\label{i-3}

\subsubsection{Validación con una organización real}
El siguiente paso será probar Economicon con usuarios, datos y una cuenta de Azure reales, comprobando permisos y separación de clientes, y medir el impacto de negocio que hoy solo es una hipótesis. Se repetirán las evaluaciones con más revisores y repeticiones, y se medirá la latencia en la infraestructura definitiva antes de aceptar cambios.

\subsubsection{Herramientas para el modelo y tiempo de respuesta}
Las consultas de propiedad, gasto y anomalías se ofrecerán al modelo como herramientas que el servidor autoriza y ejecuta, de modo que el chat pueda usarlas sin escribir SQL. Para reducir la latencia se estudiará acotar la salida sin truncarla, reutilizar contexto mediante caché y activar la recuperación híbrida tras preparar su índice. \cite{arquitectura-final}

\subsubsection{Una referencia predictiva para las anomalías}
Hoy una anomalía se detecta comparando el gasto de los últimos siete días con el de los siete anteriores: se marca si sube más de un 20\,\% y más de 100 unidades monetarias diarias. La regla es determinista y explicable, pero no distingue un crecimiento normal de un salto inesperado. Se propondrá un modelo que estime el gasto esperado de cada día: una regresión lineal como punto de partida, comparada con modelos estacionales y técnicas no supervisadas por precisión y explicabilidad. Hará falta un histórico de varias semanas o meses que hoy no existe. \cite{detector-anomalias,jup130-anomalias}

\subsubsection{Detección periódica y avisos}
Las anomalías se detectan hoy cuando la persona pulsa un botón, y cada una se guarda como un episodio por cliente, moneda y servicio con su estado de seguimiento. Se añadirá una detección periódica, apagada por defecto y coordinada con la carga de costes, con aviso primero dentro de la aplicación y después, si aporta valor, mediante una regla de Grafana. Antes de activarla se definirán su frecuencia, el tratamiento de duplicados y las condiciones de envío. \cite{detector-anomalias}

\subsubsection{Despliegue en la nube}
El ciclo desechable en AWS está preparado. Ejecutarlo permitirá medir la ruta con modelos nativos y comparar coste y latencia con la configuración actual. \cite{aws-ensayo}
